% ============================================================
% SECTION 2 — MODEL AND BACKGROUND
% ============================================================
\section{Model and Background}
\label{sec:model}

We first recall the general time-series prediction framework of
\citet{xu2023conformal} and the EnbPI algorithm, then
introduce the seasonal heteroskedasticity structure that motivates our work.

%------------------------------------------------------------
\subsection{Problem Setup}
\label{sec:model:setup}
%------------------------------------------------------------

Let $f:\R^{d}\to\R$ be an unknown function and let
$\{(X_t, Y_t)\}_{t \geq 1}$ be an observed process generated by
%
\begin{equation}
  \label{eq:dgp}
  Y_t \;=\; f(X_t) + \epsilon_t, \qquad t = 1, 2, \ldots,
\end{equation}
%
where $X_t \in \R^d$ is a feature vector (which may include lagged
values of $Y_t$ and exogenous variables) and $\epsilon_t$ is a
stochastic error term with continuous cumulative distribution function
(CDF) $F_t$. We do \emph{not} require $\epsilon_t$ to be
independent, nor that $F_t$ be the same across all $t$.

We observe the first $T$ pairs $\{(x_t, y_t)\}_{t=1}^T$ as
training data and wish to construct a prediction interval
$\hat{C}_{T+1}^{\,\alpha}$ for $Y_{T+1}$ at a user-specified
significance level $\alpha \in (0,1)$. We consider two coverage
notions of increasing strength.

\begin{definition}[Coverage guarantees]
\label{def:coverage}
A prediction interval $\hat{C}_t^{\,\alpha}$ is said to be:
\begin{itemize}
  \item \textit{conditionally valid} at time $t > T$ if
  %
  \begin{equation}
    \label{eq:cond_coverage}
    \Prob\!\bigl(Y_t \in \hat{C}_t^{\,\alpha} \,\big|\, X_t = x_t\bigr)
    \;\geq\; 1 - \alpha;
  \end{equation}

  \item \textit{marginally valid} at time $t > T$ if
  %
  \begin{equation}
    \label{eq:marg_coverage}
    \Prob\!\bigl(Y_t \in \hat{C}_t^{\,\alpha}\bigr)
    \;\geq\; 1 - \alpha.
  \end{equation}
\end{itemize}
\end{definition}

Conditional validity~\eqref{eq:cond_coverage} is substantially
stronger: it requires correct coverage for each realised feature
vector $x_t$, not just on average over the distribution of $X_t$.
\citet{barber2023conformal} show that exact finite-sample conditional
validity is generally unattainable without parametric assumptions.
Our goal, following \citet{xu2023conformal}, is to achieve
\emph{asymptotic} conditional and marginal validity.

%------------------------------------------------------------
\subsection{Oracle Prediction Interval}
\label{sec:model:oracle}
%------------------------------------------------------------

With full knowledge of $f$ and $F_t$, denote by $F_{t,Y}$ the
conditional CDF of $Y_t$ given $X_t = x_t$:
%
\[
  F_{t,Y}(y) \;=\; \Prob(Y_t \leq y \mid X_t = x_t)
             \;=\; F_t(y - f(x_t)).
\]
%
The shortest interval with exact conditional coverage $1-\alpha$ is
%
\begin{equation}
  \label{eq:oracle}
  C_t^{\,\alpha}
  \;=\;
  \bigl[f(x_t) + F_t^{-1}(\beta^*),\;
        f(x_t) + F_t^{-1}(1-\alpha+\beta^*)\bigr],
\end{equation}
%
where $F_t^{-1}(\cdot)$ is the quantile function of $F_t$ and
%
\begin{equation}
  \label{eq:betastar}
  \beta^* \;:=\;
  \operatorname*{arg\,min}_{\beta\in[0,\,\alpha]}
  \bigl[F_t^{-1}(1-\alpha+\beta) - F_t^{-1}(\beta)\bigr].
\end{equation}
%
The optimisation in~\eqref{eq:betastar} selects the lower quantile
$\beta^*$ that minimises interval width among all intervals with
coverage $1-\alpha$.  For symmetric, unimodal $F_t$ we have
$\beta^* = \alpha/2$; for skewed distributions (e.g.\ non-negative
economic variables) $\beta^* \neq \alpha/2$, justifying asymmetric
intervals.

%------------------------------------------------------------
\subsection{Background: EnbPI}
\label{sec:model:enbpi}
%------------------------------------------------------------

We briefly recall the EnbPI framework of \citet{xu2023conformal},
which our method extends.  Given a training set
$\{(x_t,y_t)\}_{t=1}^T$ and a prediction algorithm $\calA$,
EnbPI proceeds in two phases.

\paragraph{Training phase.}
EnbPI fits $B$ bootstrap estimators
$\{\hat{f}^{\,b}\}_{b=1}^B$, where $\hat{f}^{\,b} = \calA\bigl((x_i,y_i),\, i\in S_b\bigr)$
and $S_b$ is an index set drawn with replacement from $\{1,\ldots,T\}$.
For each training point $i$, the \emph{leave-one-out} (LOO)
ensemble predictor is
\[
  \hat{f}_{-i}^{\,\phi}(x_i)
  \;=\;
  \phi\!\bigl(\hat{f}^{\,b}(x_i):\, b \text{ s.t. } i \notin S_b\bigr),
\]
where $\phi$ is an aggregation function (e.g.\ sample mean).
The LOO residual at training point $i$ is
\[
  \hat{\epsilon}_i \;=\; y_i - \hat{f}_{-i}^{\,\phi}(x_i).
\]
The ordered multiset $\hat{\boldsymbol\epsilon} = \{\hat{\epsilon}_i\}_{i=1}^T$
constitutes the \emph{calibration buffer}.

\paragraph{Prediction phase.}
For each test point $x_t$, $t > T$, the point forecast is
\[
  \hat{f}_{-t}^{\,\phi}(x_t)
  \;=\;
  \phi\!\bigl(\hat{f}_{-i}^{\,\phi}(x_t):\, i = 1,\ldots,T\bigr).
\]
The optimal lower quantile level is obtained by a one-dimensional
line search:
\[
  \hat{\beta}
  \;=\;
  \operatorname*{arg\,min}_{\beta \in [0,\,\alpha]}
  \bigl[
    \Quant{1-\alpha+\beta}{\hat{\boldsymbol\epsilon}}
    -
    \Quant{\beta}{\hat{\boldsymbol\epsilon}}
  \bigr],
\]
where $Q_p(\cdot)$ denotes the $p$-th empirical quantile.
The prediction interval is
%
\begin{equation}
  \label{eq:enbpi_interval}
  \hat{C}_t^{\,\alpha}
  \;=\;
  \Bigl[
    \hat{f}_{-t}^{\,\phi}(x_t)
    + \Quant{\hat\beta}{\hat{\boldsymbol\epsilon}},\;
    \hat{f}_{-t}^{\,\phi}(x_t)
    + \Quant{1-\alpha+\hat\beta}{\hat{\boldsymbol\epsilon}}
  \Bigr].
\end{equation}
%
Every $s_0 \geq 1$ steps, once actual responses become available, the
buffer is updated by discarding the $s_0$ oldest residuals and appending
$s_0$ new ones, producing a sliding window of size $T$.

\paragraph{Coverage guarantee.}
The validity of EnbPI rests on two assumptions.

\begin{assumption}[Short-term i.i.d.\ errors; {\citealt{xu2023conformal}}]
\label{ass:xu_a1}
The errors $\{\epsilon_t\}_{t=1}^{T+1}$ are independent and
identically distributed according to a common CDF $F$, which is
Lipschitz continuous with constant $L > 0$.
\end{assumption}

\begin{assumption}[Estimation quality; {\citealt{xu2023conformal}}]
\label{ass:xu_a2}
There exists a real sequence $\{\delta_T\}_{T \geq 1}$ such that
\[
  \frac{1}{T}\sum_{t=1}^T \bigl(\hat{f}_{-t}(x_t) - f(x_t)\bigr)^2
  \;\leq\; \delta_T^2
  \qquad\text{and}\qquad
  \bigl|\hat{f}_{-(T+1)}(x_{T+1}) - f(x_{T+1})\bigr| \leq \delta_T.
\]
\end{assumption}

Under these two assumptions, \citet{xu2023conformal} establish the
following finite-sample bound on the conditional coverage gap.

\begin{theorem}[Conditional coverage gap; {\citealt[Theorem~1]{xu2023conformal}}]
\label{thm:xu_original}
Under Assumptions~\ref{ass:xu_a1} and~\ref{ass:xu_a2}, for any
$T \in \N$, $\alpha \in (0,1)$, and $\beta \in [0,\alpha]$:
\[
  \bigl|
    \Prob\!\bigl(Y_{T+1} \in \hat{C}_{T+1}^{\,\alpha}
                 \mid X_{T+1} = x_{T+1}\bigr) - (1-\alpha)
  \bigr|
  \;\leq\;
  12\sqrt{\frac{\log(16T)}{T}}
  + 4(L+1)\!\left(\delta_T^{2/3} + \delta_T\right).
\]
If $\delta_T \to 0$ as $T \to \infty$, both terms on the right-hand
side converge to zero, establishing asymptotic conditional validity.
\end{theorem}

\begin{remark}
  Assumption~\ref{ass:xu_a1} can be relaxed to errors following a
  linear process or a strongly mixing sequence, at the cost of slower
  convergence rates \citep[Corollaries~1--2]{xu2023conformal}.
  Theorem~\ref{thm:xu_original} also implies marginal validity by
  the law of total probability.
\end{remark}

%------------------------------------------------------------
\subsection{Seasonal Heteroskedasticity}
\label{sec:model:seasonal}
%------------------------------------------------------------

Assumption~\ref{ass:xu_a1} requires that all $T+1$ errors share a
common CDF $F$.  For many economic series this fails: the
distribution of forecast errors varies systematically with the
calendar season.  We now formalise this structure.

\begin{definition}[Seasonal index]
\label{def:season}
For a fixed seasonal period $\Sper \in \N$ (e.g.\ $\Sper = 12$ for
monthly data), define the season of observation $t$ as
\[
  s(t) \;=\; \bigl[(t-1) \bmod \Sper\bigr] + 1 \;\in\; \{1,\ldots,\Sper\}.
\]
The stratum of season $s$ within the training sample is
$\calT_s \;=\; \{t \in [T] : s(t) = s\}$,
with cardinality $T_s = |\calT_s|$.
\end{definition}

We adopt the following model for the error process.

\begin{assumption}[Seasonally heteroskedastic errors]
\label{ass:seasonal}
There exist constants $\sigma_1, \ldots, \sigma_\Sper > 0$ and a
CDF $G$, Lipschitz continuous with constant $L_G > 0$, such that
\begin{equation}
  \label{eq:seasonal_model}
  \epsilon_t \;=\; \sigma_{s(t)}\,\eta_t,
  \qquad
  \{\eta_t\}_{t \geq 1} \;\overset{\mathrm{i.i.d.}}{\sim}\; G, \quad
  \E[\eta_t] = 0, \quad \E[\eta_t^2] = 1.
\end{equation}
\end{assumption}

Under Assumption~\ref{ass:seasonal}, the CDF of $\epsilon_t$ depends
only on the season $s(t)$:
\begin{equation}
  \label{eq:seasonal_cdf}
  F_s(x) \;=\; G\!\left(\frac{x}{\sigma_s}\right),
  \qquad s = 1, \ldots, \Sper.
\end{equation}
Each $F_s$ is Lipschitz with constant $L_s = L_G/\sigma_s$.

\begin{remark}
Assumption~\ref{ass:seasonal} is a \emph{scale mixture} model: the
shape of the error distribution is fixed (given by $G$), and only the
scale changes across seasons.  This encompasses seasonal
ARCH effects and periodic variance models studied in
\citet{franses2004periodic}.  The restriction to a common shape $G$ is
a simplifying device; our results can be extended to the case of
season-specific shapes $G_s$ at the cost of heavier notation.
\end{remark}

\paragraph{Pooled CDF.}
If the training sample contains $T_s = T/\Sper$ observations per
season (balanced panel), the empirical CDF computed over all $T$
LOO residuals estimates the \emph{mixture} of the seasonal CDFs:
\begin{equation}
  \label{eq:pooled_cdf}
  \bar{F}(x)
  \;=\;
  \frac{1}{\Sper}\sum_{s=1}^{\Sper} F_s(x)
  \;=\;
  \frac{1}{\Sper}\sum_{s=1}^{\Sper}
      G\!\left(\frac{x}{\sigma_s}\right).
\end{equation}
When predicting in season $\sstar$, the relevant CDF is $F_{\sstar}$,
not $\bar{F}$.  We quantify the discrepancy between these two objects
via the following key quantity.

\begin{definition}[Seasonal bias]
\label{def:seasonal_bias}
For a prediction target in season $\sstar$, the \emph{seasonal bias} is
\begin{equation}
  \label{eq:seasonal_bias}
  \bsstar
  \;:=\;
  \sup_{x \in \R}\,\bigl|F_{\sstar}(x) - \bar{F}(x)\bigr|
  \;=\;
  \sup_{x \in \R}\,
  \left|
    G\!\left(\frac{x}{\sigma_{\sstar}}\right)
    -
    \frac{1}{\Sper}\sum_{s=1}^{\Sper}
         G\!\left(\frac{x}{\sigma_s}\right)
  \right|.
\end{equation}
\end{definition}

Several properties of $\bsstar$ are immediate.

\begin{lemma}[Properties of the seasonal bias]
\label{lem:bias_properties}
Under Assumption~\ref{ass:seasonal}:
\begin{enumerate}
  \item[\textup{(i)}]
    $\bsstar = 0$ if $\sigma_s = \sigma$ for all $s \in \{1,\ldots,\Sper\}$
    (seasonal homoskedasticity).

  \item[\textup{(ii)}]
    $\bsstar$ does not depend on $T$; it is a fixed property
    of the data-generating process.

  \item[\textup{(iii)}]
    If $G$ is strictly increasing and season $\sstar$ has the
    strictly largest or the strictly smallest scale (that is,
    $\sigma_{\sstar} = \max_s \sigma_s$ or
    $\sigma_{\sstar} = \min_s \sigma_s$, with
    $\sigma_s \neq \sigma_{\sstar}$ for at least one $s$),
    then $\bsstar > 0$.
\end{enumerate}
\end{lemma}

\begin{proof}
  Part~(i): if all scales are equal, then $F_s(x) = G(x/\sigma)$ for
  every $s$, so $F_{\sstar} = \bar{F}$ pointwise and $\bsstar = 0$.
  Part~(ii) follows directly from Definition~\ref{def:seasonal_bias},
  since neither $F_{\sstar}$ nor $\bar{F}$ depend on the sample size $T$.
  Part~(iii): suppose $\sigma_{\sstar} = \max_s \sigma_s$ and
  $\sigma_{s'} < \sigma_{\sstar}$ for some $s'$.  Fix any $x > 0$.
  Then $x/\sigma_{\sstar} \leq x/\sigma_s$ for every $s$, with strict
  inequality for $s = s'$; since $G$ is strictly increasing,
  $G(x/\sigma_{\sstar}) \leq G(x/\sigma_s)$ with strict inequality for
  $s'$, whence
  $\bar{F}(x) > G(x/\sigma_{\sstar}) = F_{\sstar}(x)$ and therefore
  $\bsstar \geq \bar{F}(x) - F_{\sstar}(x) > 0$.  The case
  $\sigma_{\sstar} = \min_s \sigma_s$ is symmetric.
\end{proof}

Part~(ii) of Lemma~\ref{lem:bias_properties} is the key observation:
$\bsstar$ is \emph{independent of the sample size $T$}.  Any
method that calibrates its prediction intervals using the pooled
distribution $\bar{F}$ (rather than the season-specific $F_{\sstar}$)
will therefore incur a coverage error bounded away from zero,
regardless of how much data is available.  Section~\ref{sec:pooled_gap}
makes this precise.

\paragraph{Example: Gaussian errors.}
Suppose $G = \Phi$ (standard normal), $\Sper = 12$, with
$\sigma_H = 2$ for $k = 3$ high-volatility seasons (set $\calH$)
and $\sigma_L = 1$ for the remaining nine.  Then:
\[
  \bar{F}(x)
  \;=\;
  \tfrac{3}{12}\,\Phi\!\left(\tfrac{x}{2}\right)
  +
  \tfrac{9}{12}\,\Phi(x),
\]
and for a prediction in a high-volatility season $\sstar \in \calH$:
\[
  b_H
  \;=\;
  \sup_x\!
  \left|
    \Phi\!\left(\tfrac{x}{2}\right)
    -\tfrac{1}{4}\Phi\!\left(\tfrac{x}{2}\right)
    -\tfrac{3}{4}\Phi(x)
  \right|
  \;=\;
  \tfrac{3}{4}\sup_x\!\left|\Phi\!\left(\tfrac{x}{2}\right) - \Phi(x)\right|.
\]
Setting the derivative $\frac{d}{dx}\bigl[\Phi(x/2)-\Phi(x)\bigr] = 0$
yields the critical point $x^* = -\sqrt{8\ln 2/3} \approx -1.36$,
at which $\Phi(x^*/2) \approx 0.248$ and $\Phi(x^*) \approx 0.087$.
Hence:
\[
  b_H \;=\; \tfrac{3}{4}(0.248 - 0.087) \;\approx\; 0.12.
\]

Since the coverage gap bound is $2b_H \approx 0.24$, pooled
calibration can introduce a \emph{persistent} coverage gap of
up to $24$ percentage points in high-volatility seasons at a
nominal level of $\alpha = 0.10$.
Crucially, $b_H$ is independent of $T$ by
Lemma~\ref{lem:bias_properties}(ii): collecting more data does not
reduce it.  In contrast, the statistical fluctuation in the empirical
CDF decays at rate $O(T^{-1/2})$.  Therefore, for all sufficiently
large $T$, the dominant source of coverage error in high-volatility
seasons is the fixed seasonal bias $2b_H$, not sampling variability.
Proposition~\ref{prop:pooled_gap} makes this precise.\medskip

We close this section by noting that model~\eqref{eq:seasonal_model}
nests the homoskedastic case of \citet{xu2023conformal}
(set $\sigma_s = \sigma$ for all $s$) and reduces to their Assumption~1
when $\Sper = 1$.  All results of \citet{xu2023conformal} therefore
hold as special cases of our framework.
